% Dibangkitkan oleh revisi/hasil/gerbang_batch/bab5_jalan.py. Jangan diubah dengan tangan.
\begin{table}[htbp]
\caption{Celah ARLS-PPO ke Sequential Lookahead: $\lambda$ tanpa label dikurangi acuan (selang 95\%)}
\label{tab:uji-butir1}
\centering\footnotesize
\setlength{\tabcolsep}{4pt}
\begin{adjustbox}{max width=\textwidth}
\begin{tabular}{p{6cm}ll}
\toprule
Pasangan $\lambda$ & Uk.~1 (poin) & Uk.~2 (poin) \\
\midrule
0,1 lawan 0,3 & $+9{,}14$ [$+2{,}07$; $+17{,}75$]$^{*}$ & $-0{,}62$ [$-2{,}20$; $+1{,}05$] \\
1 lawan interpolasi log 0,3 dan 3 & $-6{,}23$ [$-12{,}29$; $-1{,}64$]$^{*}$ & $-1{,}13$ [$-2{,}68$; $+0{,}38$] \\
10 lawan interpolasi log 3 dan 30 & $+5{,}14$ [$+3{,}64$; $+6{,}68$]$^{*}$ & $+4{,}31$ [$+3{,}18$; $+5{,}47$]$^{*}$ \\
pelengkap: rata-rata tanpa label lawan rata-rata berlabel & $+7{,}36$ [$+4{,}68$; $+10{,}51$]$^{*}$ & $+1{,}98$ [$+1{,}17$; $+2{,}92$]$^{*}$ \\
\bottomrule
\end{tabular}
\end{adjustbox}
\par\smallskip\footnotesize\textit{Catatan:} Positif berarti celah ke Sequential Lookahead di $\lambda$ tanpa label lebih besar daripada di acuan berlabel. Batas atas selang adalah $X$ pada kalimat klaim butir~1.
\end{table}
